\documentclass[12pt, a4paper]{report}

% 

\usepackage{elegant-cours}
\makeindex

\begin{document}

\MaPageDeGarde{Chapitre XIV-B : Applications Linéaires}{Retranscrit par Samy Youssoufine}{./assets/logo.png}{WIP. Peut contenir des erreurs ou des sections incomplètes.}

\tableofcontents
\clearpage

\vspace*{\fill}
Ce chapitre est consacré aux applications linéaires, qui sont des fonctions entre espaces vectoriels qui préservent les opérations de l'addition et de la multiplication par un scalaire.
\vspace*{\fill}

\chapter{Généralités}

\section{Définitions et propriétés}

\begin{definition}
    Soient $E$ et $F$ deux $\Kset$-ev. et $f : E \to F$.
    \begin{enumerate}
        \item $f$ est dite une application linéaire de $E$ vers $F$ lorsque $\forall x, y \in E, \forall \lambda \in \Kset$,
        \[f(x + y) = f(x) + f(y) \quad \text{et} \quad f(\lambda x) = \lambda f(x).\]
        On note $\mathcal{L}(E, F)$ l'ensemble des applications linéaires de $E$ vers $F$.
        \item Si $f \in \mathcal L(E, F)$ est bijective, on dit que $f$ est un isomorphisme d'espaces vectoriels de $E$ vers $F$.\\On note $\Isom(E, F)$ l'ensemble des isomorphismes de $E$ vers $F$.
        \item Si $f \in \mathcal L(E,E)$, noté $\mathcal L(E)$, alors $f$ est dite un endomorphisme de $E$.
        \item Si $f\in\mathcal L(E)$ est bijective, alors $f$ est dite un automorphisme de $E$.\\On note $\Aut(E)$ l'ensemble des automorphismes de $E$.
    \end{enumerate}
\end{definition}

\begin{example}
    \begin{outline}[enumerate]
        \1 Soient $a_0,\dots,a_n\in\Kset$ et $\varphi:\begin{matrix}
            \Kset_n[X] &\to &\Kset^{n+1}\\
            P &\mapsto &(P(a_0),\dots,P(a_n))
        \end{matrix}$, alors $\varphi\in\mathcal L(\Kset_n[X],\Kset^{n+1})$.
        \2 En effet, $\Kset_n[X]$ et $\Kset^{n+1}$ sont des $\Kset$-ev. et pour tous $P,Q\in\Kset_n[X]$ et $\lambda\in\Kset$, on a
        \begin{align*}
            \varphi(P+\lambda Q) &= ((P+\lambda Q)(a_0),\dots,(P+\lambda Q)(a_n))\\
            &= (P(a_0)+\lambda Q(a_0),\dots,P(a_n)+\lambda Q(a_n))\\
            &= (P(a_0),\dots,P(a_n)) + \lambda(Q(a_0),\dots,Q(a_n))\\
            &= \varphi(P) + \lambda\varphi(Q).
        \end{align*}
        \2 Cette application est un isomorphisme si et seulement si les $a_i$ sont deux à deux distincts. On peut aussi montrer qu'elle est surjective en utilisant les polynômes de Lagrange (interpolation lagrangienne)\index{polynômes de Lagrange}.
        \1 Soient $a < b \in \R$. On pose $\varphi : \begin{matrix}
            \mathcal C([a,b],\R) &\to &\R\\
            f &\mapsto &\int_a^b f
        \end{matrix}$, alors $\varphi\in\mathcal L(\mathcal C([a,b],\R),\R)$.
        \1 $\Delta : \begin{matrix}
            \Kset[X] &\to &\Kset[X]\\
            P &\mapsto &P(X+1) - P(X)
        \end{matrix}$ est une application linéaire de $\Kset[X]$ vers lui-même (endomorphisme de $\Kset[X]$).
    \end{outline}
\end{example}

\begin{property}
    Soient $E,F$ deux $\Kset$-ev. et $f \in \mathcal L(E,F)$.
    \begin{enumerate}
        \item $f(0_E) = 0_F$.
        \item $\forall x_1,\dots,x_n\in E, \forall \lambda_1,\dots,\lambda_n\in\Kset$, $f(\lambda_1 x_1 + \dots + \lambda_n x_n) = \lambda_1 f(x_1) + \dots + \lambda_n f(x_n)$.
    \end{enumerate}
\end{property}

\begin{proof}
    \begin{outline}[enumerate]
        \1 $f(0_E) = f(0_E + 0_E) = f(0_E) + f(0_E)$, donc $f(0_E) = 0_F$.
        \1 Par récurrence sur $n$. Le cas $n=1$ est trivial. Supposons que la propriété soit vérifiée pour $n$, montrons-la pour $n+1$. Soient $x_1,\dots,x_{n+1}\in E$ et $\lambda_1,\dots,\lambda_{n+1}\in\Kset$, alors
        \begin{align*}
            f(\lambda_1 x_1 + \dots + \lambda_{n+1} x_{n+1}) &= f((\lambda_1 x_1 + \dots + \lambda_n x_n) + \lambda_{n+1} x_{n+1})\\
            &= f(\lambda_1 x_1 + \dots + \lambda_n x_n) + f(\lambda_{n+1} x_{n+1})\\
            &= (\lambda_1 f(x_1) + \dots + \lambda_n f(x_n)) + \lambda_{n+1} f(x_{n+1})\\
            &= \lambda_1 f(x_1) + \dots + \lambda_{n+1} f(x_{n+1}) \text{ par H.R.}
        \end{align*}
    \end{outline}
\end{proof}

\begin{property}
    Soient $E,F$ deux $\Kset$-ev.
    $$\parens{\mathcal L(E,F), +, \cdot} \text{ est un } \Kset\text{-ev.}$$
\end{property}

\begin{proof}
    \begin{outline}
        \1 On a $\mathcal L(E,F) \subset F^E$.
        \1 L'application nulle (notée $0_{F^E}$) est bien linéaire. Donc $0_F \in \mathcal L(E,F)$.
        \1 Soient $f,g \in \mathcal L(E,F)$ et $\alpha \in \Kset$. Montrons que $f + \alpha g \in \mathcal L(E,F)$. Soient $x,y \in E$ et $\lambda \in \Kset$, alors :
        \begin{align*}
            (f + \alpha g)(x + \lambda y) &= f(x + \lambda y) + \alpha g(x + \lambda y)\\
            &= (f(x) + \lambda f(y)) + \alpha(g(x) + \lambda g(y))\\
            &= (f(x) + \alpha g(x)) + (\lambda f(y) + \alpha \lambda g(y))\\
            &= (f + \alpha g)(x) + \lambda (f + \alpha g)(y)
        \end{align*}
        \1 Ainsi, $f + \alpha g \in \mathcal L(E,F)$.
        \1 Par conséquent, $\mathcal L(E,F)$ est stable par addition et multiplication par un scalaire, et contient l'élément neutre pour l'addition. Donc $\parens{\mathcal L(E,F), +, \cdot}$ est un sous-espace vectoriel de $F^E$. Or $F^E$ est un $\Kset$-ev. (chapitre précédent), donc $\mathcal L(E,F)$ est un $\Kset$-ev.
    \end{outline}
\end{proof}

\section{L'anneau $\mathcal L(E)$}

\subsection{Définitions et propriétés de base}

\begin{property}
    Soient $E,F,G$ trois $\Kset$-ev. et $f \in \mathcal L(E,F)$, $g \in \mathcal L(F,G)$, alors $g \circ f \in \mathcal L(E,G)$.
\end{property}

\begin{proof}
    Soient $x,y \in E$ et $\lambda \in \Kset$, alors :
    \begin{align*}
        (g \circ f)(x + \lambda y) &= g(f(x + \lambda y))\\
        &= g(f(x) + \lambda f(y))\\
        &= g(f(x)) + \lambda g(f(y))\\
        &= (g \circ f)(x) + \lambda (g \circ f)(y)
    \end{align*}
    Ainsi, $g \circ f$ est une application linéaire de $E$ vers $G$, c'est-à-dire $g \circ f \in \mathcal L(E,G)$.
\end{proof}

\begin{property}
    $$(\mathcal L(E), +, \circ) \text{ est un anneau non-commutatif.}$$
\end{property}

\begin{proof}
    On doit montrer que :
    \begin{outline}
        \1 La loi $\circ$ est une loi de composition interne sur $\mathcal L(E)$.
        \1 $(\mathcal L(E), +)$ est un groupe abélien.
        \1 Il existe un élément neutre pour $\circ$ dans $\mathcal L(E)$.
        \1 $\circ$ est associative.
        \1 $\circ$ est distributive à gauche et à droite par rapport à $+$.
    \end{outline}
    On procède alors à la vérification de chacune de ces propriétés :
    \begin{outline}
        \1 On a $(\mathcal L(E), +, \cdot)$ est un $\Kset$-ev. Donc $(\mathcal L(E), +)$ est un groupe abélien.
        \1 D'après la propriété précédente, on a $\circ$ est une loi de composition interne dans $\mathcal L(E)$.
        \1 On a $\Id_E \in \mathcal L(E)$ et pour tout $f \in \mathcal L(E)$, on a $\Id_E \circ f = f \circ \Id_E = f$. Donc $\Id_E$ est l'élément neutre pour la loi $\circ$.
        \1 On a $\circ$ est associative, c'est-à-dire que pour tous $f,g,h \in \mathcal L(E)$, on a $(f \circ g) \circ h = f \circ (g \circ h)$.
        \1 Aussi, $\circ$ est distributive par rapport à $+$, c'est-à-dire que pour tous $f,g,h \in \mathcal L(E)$, on a $f \circ (g + h) = f \circ g + f \circ h$ et $(f + g) \circ h = f \circ h + g \circ h$.
        \1 Cependant, $\circ$ n'est pas commutative en général, c'est-à-dire qu'il existe des $f,g \in \mathcal L(E)$ tels que $f \circ g \neq g \circ f$ (on démontre ça par contre-exemple).
            \2 On prend $E = \R^n$ avec $n \ge 2$
            \2 On pose $f : \begin{matrix}
            R^n &\to &R^n\\
            (x_1,\dots,x_n) &\mapsto &(x_1,0,\dots,0)
            \end{matrix}$ et $g : \begin{matrix}
            R^n &\to &R^n\\
            (x_1,\dots,x_n) &\mapsto &(x_2,0,\dots,0)
            \end{matrix}$
            \2 Alors $\forall (x_1,\dots,x_n) \in \R^n$, on a $f(x) = x_1\cdot e_1$ et $g(x) = x_2\cdot e_1$.
            \2 $(f \circ g)(e_2) = f(g(e_2)) = f(e_1) = e_1$ alors que $(g \circ f)(e_2) = g(f(e_2)) = g(0) = 0_{\R^n}$. Donc $f \circ g \neq g \circ f$.
            \2 Ainsi, $\circ$ n'est pas commutative en général.
            \3 On pourra aussi noter que cet anneau n'est ni intègre (car il n'est pas commutatif et un anneau intègre est défini comme étant un anneau commutatif sans diviseur de zéro) ni semi-intègre\index{Semi-intègre} (car, dans ce dernier contre-exemple $f\circ g(x) = 0_{R^n}$ mais $f \neq 0$ et $g \neq 0$).
    \end{outline}
\end{proof}

\begin{remark}
    Comme la loi multiplicative de $\mathcal L(E)$ est $\circ$, on utilise alors les notations $f\cdot g$ ou $fg$ pour désigner $f \circ g$.

    Donc $fg(x) = f(g(x))$ pour tous $f,g \in \mathcal L(E)$ et $x \in E$, et non pas $f(x)g(x)$, sachant que nous n'avons même pas défini la loi multiplicative dans $E$.

    On notera aussi, pour tout $f \in \mathcal L(E)$ et $n\ge 1$, $f^n = \underbrace{f \circ f \circ \dots \circ f}_{n\text{ fois}}$ et $f^0 = \Id_E$.

    \label{remarque:notation-multiplication-anneau-LE}
\end{remark}

\subsection{Groupes linéaires}

\begin{definition}[Groupe linéaire]
    Soit $E$ un $\Kset$-ev.

    L'ensemble des éléments inversibles de $(\mathcal L(E), +, \circ)$ est un groupe non-commutatif, appelé le groupe linéaire de $E$ et noté $\GL(E)$.
\end{definition}

\begin{example}
    On pose : $f : \begin{matrix}
        \R^2 &\to &\R^2\\
        (x,y) &\mapsto &(x-y,x+2y)
    \end{matrix}$

    On a $f \in \mathcal L(\R^2)$.

    Soit $(a,b) \in \R^2$. On résout $f(x,y) = (a,b)$, c'est-à-dire le système d'équations linéaires suivant :

    \[\begin{cases}
    x - y = a\\
    x + 2y = b
    \end{cases}\]

    Après résolutions, on trouve $x = \frac{2a + b}{3}$ et $y = \frac{b - a}{3}$. La surjectivité est donc démontrée, et l'unicité est aussi démontrée implicitement (vu qu'on a trouvé une unique expression). Donc $f$ est bijective, et par conséquent $f \in \GL(\R^2)$, et $f^{-1} : \begin{matrix}
        \R^2 &\to &\R^2\\
        (a,b) &\mapsto &\left(\frac{2a + b}{3}, \frac{b - a}{3}\right)
    \end{matrix}$.

    \textit{On verra dans le chapitre suivant comment utiliser les matrices pour démontrer l'inversibilité d'une application linéaire, et pour calculer son inverse. On utilisera une matrice associée canoniquement à $f$ et on montrera que cette dernière est inversible, et que son inverse est la matrice associée canoniquement à $f^{-1}$.}
\end{example}

\begin{remark}[Polynôme d'endomorphisme (Programme Spé.)]
    Soient $f$ un endomorphisme de $E$, et $P = \sum_{k=0}^n a_k X^k \in \Kset[X]$.

    On appelle $P(f) = \sum_{k=0}^n a_k f^k = a_0\cdot\Id_E + a_1\cdot f + \dots + a_n\cdot f^n$ le polynôme d'endomorphisme de $f$ associé à $P$.

    \textit{(On rappelle que la notation $f^k$ est définie dans la remarque précédente \ref{remarque:notation-multiplication-anneau-LE}...)}

    \begin{enumerate}
        \item Si $P(f) = 0_{\mathcal L(E)}$ (endomorphisme nul), alors $P$ est appelé un \textbf{polynôme annulateur} de $f$.
        \item Si $P(f) = 0_{\mathcal L(E)}$ avec $n=\deg(P) \ge 1$ et $a_0 \ne 0_\Kset$, alors $a_1\cdot f + \dots + a_n\cdot f^n = -a_0\cdot\Id_E$. Comme $f$ est linéaire, on peut composer pour obtenir $f\circ\left(a_1\cdot\Id_E + \dots + a_n\cdot f^{n-1}\right)= -a_0\cdot\Id_E$.\\Donc, comme $a_0 \ne 0$, on a $f\circ\left((-a_0^{-1})(a_1\cdot\Id_E + \dots + a_n\cdot f^{n-1})\right) = \Id_E$. Donc $f \in \GL(E)$ est inversible, et $f^{-1} = (-a_0^{-1})(a_1\cdot\Id_E + \dots + a_n\cdot f^{n-1})$.
        \item Soit $f \in \mathcal L(E)$. On pose : $\Kset[f] \overset{\text{noté}}{=} \{P(f) : P \in \Kset[X]\}$. Alors $(\Kset[f],+,\circ)$ est un anneau commutatif (car il est stable par addition et par composition).
    \end{enumerate}

    \textit{Note additionnelle (Théorème de Cayley-Hamilton) : Tout endomorphisme $f$ de $E$ admet un polynôme annulateur, et même plus précisément, le polynôme caractéristique de $f$ est un polynôme annulateur de $f$.}
\end{remark}
\todo{recop}

\subsection{Endomorphismes nilpotents}

\begin{definition}[Endomorphisme nilpotent]
    Soient $E$ un $\Kset$-ev. et $f \in \mathcal L(E)$. On dit que $f$ est nilpotent lorsque $\exists k \in \N^*$ tel que $f^k = 0_{\mathcal L(E)}$.

    On appelle $p=\min(\{k \in \N^* : f^k = 0_{\mathcal L(E)}\})$ l'indice de nilpotence de $f$.

    On a donc $f^p = 0_{\mathcal L(E)}$, $f^{p-1} \neq 0_{\mathcal L(E)}$ et $\forall k \ge p \in \N^*, f^k = 0_{\mathcal L(E)}$. 
\end{definition}

\begin{example}
    On pose $D : \begin{matrix}
        \Kset_n[X] &\to &\Kset_n[X]\\
        P &\mapsto &P'
    \end{matrix}$.

    On a :
    \begin{outline}
        \1 $D \in \mathcal L(\Kset_n[X])$
        \1 $\forall k \ge 1, \forall P \in \Kset_n[X], D^k(P) = P^{(k)}$ (la $k$-ième dérivée de $P$).
    \end{outline}

    Donc $\forall P \in \Kset_n[X], D^{n+1}(P) = P^{(n+1)} = 0$.

    Et on a $D^n(X^n) = n! \neq 0$. Donc $D$ est nilpotent d'indice de nilpotence $n+1$.
\end{example}

\begin{exercise}
    Soit $f \in \mathcal L(E)$ un endomorphisme nilpotent. Montrer que $\Id_E -f \in \GL(E)$.
    
    \textit{(Analogie avec les éléments nilpotents d'un anneau commutatif, où $1 - x$ est inversible pour tout élément nilpotent $x$)}

    \textbf{Solution :}

    \begin{outline}
        \1 $f$ est nilpotent, donc $\exists p \in \N^*$ tel que $f^p = 0_{\mathcal L(E)}$.

        \1 On a $\Id_E = \Id_E^p - f^p = (\Id_E - f)\circ\left(\sum_{k=0}^{p-1} \Id_E^k \circ f^{p-1-k}\right)$.

        \1 Donc $\Id_E - f$ est inversible, et son inverse est $\sum_{k=0}^{p-1} \Id_E^k \circ f^{p-1-k}$.
    \end{outline}
\end{exercise}

\begin{remark}
    Comme $\mathcal L(E)$ est un anneau, alors $\forall f,g \in \mathcal L(E)$ tels que $fg = gf$, on a $\forall n \in \N, (f+g)^n = \sum_{k=0}^n \legbinom{n}{k} f^k g^{n-k}$ (binôme de Newton).
\end{remark}

\begin{exercise}
    On pose : $u : \begin{matrix}
    \R^2 &\to &\R^2\\
    (x,y) &\mapsto &(x+2y,y)
    \end{matrix}$.

    Calculer $u^n$ pour tout $n \in \N^*$, puis déterminer un polynôme annulateur de cet endomorphisme.

    \textbf{Solution :}

    \begin{outline}
        \1 On voit clairement que $u^n : (x,y) \mapsto (x + 2ny, y)$. Ce résultat est démontré par récurrence sur $n$.
        \1 On pose $u = \Id_{\R^2} + v$ avec $v : (x,y) \mapsto (2y,0)$.
        \2 Le cas $n=1$ est trivial.
        \2 On a $u^n = \sum_{k=0}^{n}\legbinom{n}{k} \Id_{\R^2}^k \circ v^{n-k}$ sachant que $v\circ \Id_{\R^2} = \Id_{\R^2}\circ v$.
        \2 Or $v^2 = 0_{\mathcal L(\R^2)}$, donc $u^n = \Id_{\R^2} + n\cdot v$.
        \2 Donc $u^n = \Id_{\R^2} + n\cdot v$.
        \2 D'où $u^n : (x,y) \mapsto (x + 2ny, y)$.
        \1 Pour déterminer un polynôme annulateur de $u$, on repart de l'égalité $u^n = \Id_{\R^2} + n\cdot v$.
        \2 On a $u^n = \Id_{\R^2} + n\cdot (u - \Id_{\R^2})$.
        \2 Donc $u^n -nu + (n-1)\Id_{\R^2} = 0_{\mathcal L(\R^2)}$.
        \2 Ainsi, un polynôme annulateur de $u$ est $P(X) = X^n - nX + (n-1)$.
        \1 \textbf{Application utile des polynômes d'endomorphisme :} On peut aussi utiliser les polynômes d'endomorphisme pour calculer $u^n$ à partir de ce polynôme annulateur, en effectuant la division euclidienne\index{Division euclidienne} de $X^n$ par $(X-1)^2$/
    \end{outline}

\end{exercise}

\chapter{Image et noyau d'une application linéaire}

\section{Définitions et propriétés}

\begin{property}
    Soient $E,F$ deux $\Kset$-ev. et $f \in \mathcal L(E,F)$.
    \begin{enumerate}
        \item Si $H$ est un sev de $E$ alors $f(H)$ est un sev de $F$.
        \item Si $H$ est un sev de $F$ alors $f^{-1}(H)$ est un sev de $E$.
    \end{enumerate}
\end{property}

\begin{proof}
    recop <З
\end{proof}

\begin{definition}
    Soient $E,F$ deux $\Kset$-ev. et $f \in \mathcal L(E,F)$.
    \begin{enumerate}
        \item L'image de $f$ est l'ensemble $\Imf(f) = f(E) = \{f(x) : x \in E\} \subset F$.
        \item Le noyau de $f$ est l'ensemble $\Ker(f) = f^{-1}(\{0_F\}) = \{x \in E : f(x) = 0_F\} \subset E$.
    \end{enumerate}
\end{definition}

\begin{property}
    Soit $f \in \mathcal L(E,F)$.
    \begin{enumerate}
        \item $f$ est injective $\iff$ $\Ker(f) = \{0_E\}$.
        \item $f$ est surjective $\iff$ $f(E) = F$.
    \end{enumerate}
\end{property}

\begin{proof}
    La preuve est très similaire à celle faite dans le chapitre 10, propriété 17, pages 29-30. On utilise juste les lois $+$ qu'on vient de définir maintenant au lieu des lois $*$ et $T$ définies dans le chapitre 10.
\end{proof}

\begin{exercise}
    \begin{outline}
        \1 \textbf{Exercice 1 :} Soit $f: \begin{matrix}
        \R^3 &\to &\R^2\\
        (x,y,z) &\mapsto &(x+y-z, 2x+y+z)
        \end{matrix}$.
        \2 Montrer que $f \in \mathcal L(\R^3,\R^2)$.
        \2 Déterminer $\Ker(f)$ et $\Imf(f)$.
        \1 \textbf{Exercice 2 :} Soit $\Delta : \begin{matrix}
        \Kset_n[X] &\to &\Kset_n[X]\\
        P &\mapsto &P(X+1) - P(X)
        \end{matrix}$.
        \2 Mêmes deux premières questions que pour $f$.
        \1 \textbf{Exercice 3 :} Soient $a_0,\dots,a_n \in \Kset$ deux à deux distincts, et $\varphi : \begin{matrix}
        \Kset_n[X] &\to &\Kset^{n+1}\\
        P &\mapsto &(P(a_0),\dots,P(a_n))
        \end{matrix}$.
        \2 Montrer que $\varphi \in \Isom(\Kset_n[X],\Kset^{n+1})$.
        \1 \textbf{Exercice 4 :} Soit $f \in \mathcal L(E,F)$ et $B$ une base de $E$ (pas forcément finie).
        \2 Montrer que $f$ est injective $\iff$ $f(B)$ est libre.
        \2 Montrer que $f$ est surjective $\iff$ $f(B)$ est génératrice de $F$.
        \2 Montrer que $f \in \Isom(E,F) \iff f(B)$ est une base de $F$.
    \end{outline}
\end{exercise}

\textbf{Solutions :}

\textit{Exercice 1 :}

\begin{outline}
    \1 La démonstration de la linéarité de $f$ est triviale et se fait en utilisant la définition de base.
    \1 $\Ker(f) = \{(x,y,z) \in \R^3 : x+y-z = 0 \text{ et } 2x+y+z = 0\}$. En résolvant ce système d'équations linéaires, on trouve que $\Ker(f) = \{(-2t,3t,t) : t \in \R\}=\R\cdot v$ avec $v = (-2,3,1)$. On procède par équivalences pour démontrer implicitement l'égalité.
    \1 L'image de $f$ est l'ensemble des éléments de $\R^2$ de la forme $(x+y-z, 2x+y+z)$ pour $(x,y,z) \in \R^3$. On peut poser deux vecteurs $u_1 = (1,2) = f(e_1)$ et $u_2 = (1,1) = f(e_2)$, et on voit que $\Vect(u_1,u_2) \subset \Imf(f) \subset \R^2$, or $(u_1,u_2)$ est libre alors $\dim(\Vect(u_1,u_2)) = 2$, donc $\Vect(u_1,u_2) = \R^2$. Donc $\Imf(f) = \R^2$. $f$ est donc surjective.
\end{outline}

\textit{Exercice 2 :} On suit la même démarche... (pas corrigé en classe).

\textit{Exercice 3 :}

\begin{outline}
    \1 L'injectivité de $\varphi$ est démontrée en utilisant le fait que les $a_i$ sont deux à deux distincts, donc le polynôme admettra $n+1$ racines distinctes, et par conséquent, il sera nul.
    \1 La surjectivité de $\varphi$ est démontrée en utilisant les polynômes de Lagrange (interpolation lagrangienne)\index{polynômes de Lagrange}.
    \2 On pose $l_k = \prod_{\substack{i=0 \\ i \ne k}}^n \frac{X - a_i}{a_k - a_i}$ pour $k = 0,\dots,n$. Alors $l_k(a_j) = \delta_{jk}$ (le symbole de Kronecker), c'est-à-dire que $l_k(a_j) = 1$ si $j=k$ et $l_k(a_j) = 0$ sinon.
    \2 Soit $(y_0,\dots,y_n) \in \Kset$, alors on pose $P = \sum_{k=0}^n y_k l_k$. On a $P(a_j) = \sum_{k=0}^n y_k l_k(a_j) = y_j$ pour $j=0,\dots,n$. Donc $\varphi(P) = (P(a_0),\dots,P(a_n)) = (y_0,\dots,y_n)$, et par conséquent, $\varphi$ est surjective.
    \1 On en déduit que $\varphi$ est un isomorphisme de $\Kset_n[X]$ vers $\Kset^{n+1}$, et par conséquent, $\varphi \in \Isom(\Kset_n[X],\Kset^{n+1})$.
    \1 \textit{Note additionnelle :} On peut aussi déterminer la bijection réciproque de $\varphi$. On a donc $\varphi^{-1} : \begin{matrix}
\Kset^{n+1} &\to &\Kset_n[X]\\
(y_0,\dots,y_n) &\mapsto &\sum_{k=0}^n y_k l_k
\end{matrix}$.
    \2 Et on voit que $\sum_{k=0}^n l_k = \varphi^{-1}(1,\dots,1) = \varphi^{-1}(\varphi(1)) = 1$. Donc $l_0 + \dots + l_n = 1$.
    \2 C'est donc un résultat assez puissant : la somme des polynômes élémentaires de Lagrange est égale à 1.
\end{outline}

Le reste des exercices est laissé en tant que preuve des propriétés suivantes.

\begin{property}
    Soient $f \in \mathcal L(E,F)$ et $B$ une base de $E$ (pas forcément finie).

    \begin{enumerate}
        \item $f$ est injective $\iff$ $f(B)$ est libre.
        \item $f$ est surjective $\iff$ $f(B)$ est génératrice de $F$.
        \item $f \in \Isom(E,F) \iff f(B)$ est une base de $F$.
    \end{enumerate}
\end{property}

\begin{proof}
    \begin{outline}
        \1 Démonstration de la première équivalence :
        \2 ($\implies$) Dans le sens direct :
        \3 On a $f(B) = \{f(e) \tq e \in B\}$.
        \3 Soient $u_1,\dots,u_n \in f(B)$ avec $n \in \N^*$. Alors $\exists e_1,\dots,e_n \in B$ tels que $u_i = f(e_i)$ pour $i=1,\dots,n$.
        \3 Soient $\lambda_1,\dots,\lambda_n \in \Kset$ tels que $\sum_{k=1}^{n}\lambda_k u_k = 0_F$.
        \begin{align*}
            \implies \sum_{k=1}^{n}\lambda_k u_k = 0_F &\implies \sum_{k=1}^{n}\lambda_k f(e_k) = 0_F\\
            &\implies f\left(\sum_{k=1}^{n}\lambda_k e_k\right) = 0_F\\
            &\implies \sum_{k=1}^{n}\lambda_k e_k \in \Ker(f)\\
            &\implies \sum_{k=1}^{n}\lambda_k e_k = 0_E \text{ (car $f$ est injective)}\\
            &\implies \lambda_1 = \dots = \lambda_n = 0
        \end{align*}
        \3 Donc $(u_1,\dots,u_n)$ est libre, et par conséquent, $f(B)$ est libre.
        \2 ($\impliedby$) Dans le sens inverse :
        \3 Soit $x \in \Ker(f)$, alors $x\in E$ et $f(x) = 0_F$.
        \3 Comme $x \in E$ et $B$ est une base de $E$, alors $\exists e_1,\dots,e_n \in B$ et $\lambda_1,\dots,\lambda_n \in \Kset$ tels que $x = \sum_{k=1}^{n}\lambda_k e_k$.
    \end{outline}
\end{proof}
\todo{recop}

\begin{exercise}
    Soit $n \in \N$. Montrer que $\exists! P \in \R_{n+1}[X]$ tel que $P(0) = 0$ et $P(X+1) - P(X) = X^n$.

    \textbf{Solution :}

    On pose $E = \{P \in \R_{n+1}[X] : P(0) = 0\}$.

    $E$ est un sev de $\R_{n+1}[X]$ (car $E$ est non vide, et stable par addition et multiplication par un scalaire).

    On pose $f: \begin{matrix}
    E &\to &\R_n[X]\\
    P &\mapsto &P(X+1) - P(X)
    \end{matrix}$.

    \begin{outline}
        \1 \textbf{L'application est bien définie} parce que $\forall P \in E$, $\deg(P) \le n+1$ donc $\deg(P(X+1) - P(X)) \le n$ (coefficients dominants s'annulent).

        \1 Donc $f(P) \in \R_n[X]$ pour tout $P \in E$, et par conséquent, $f$ est une application de $E$ vers $\R_n[X]$.

        \1 $f$ est linéaire (démontré facilement).

        \1 On pose $\forall P \in E, P = \sum_{k=1}^n a_k X^k$ (le terme constant est nul car $P(0) = 0$), alors $f(P) = \sum_{k=1} a_k ((X+1)^k - X^k)$.

        \1 Or $B = (X,\dots,X^{n+1})$ est de degré échelonné, donc $B$ est libre, et par conséquent, $B$ est une base de $E$.

        \1 On a $f(B) = \{(X+1) - X, (X+1)^2 - X^2, \dots, (X+1)^{n+1} - X^{n+1}\}$. Donc $f(B)$ est de degré échelonné, et par conséquent, $f(B)$ est libre, et comme $\card{f(B)} = \card B = n+1 = \dim(\R_n[X])$, alors $f(B)$ est une base de $\R_n[X]$.

        \1 Alors $f$ est un isomorphisme de $E$ vers $\R_n[X]$, et par conséquent, $f$ est bijective. Donc $\exists! P \in E$ tel que $f(P) = X^n$, c'est-à-dire tel que $P(X+1) - P(X) = X^n$ et $P(0) = 0$.
    \end{outline}
\end{exercise}

\begin{application}
    Calculer $\sum_{k=0}^{n}k^3$ pour $n \in \N^*$.

    \textbf{Solution :}

    \begin{align*}
        f(X^4) &= (X+1)^4 - X^4\\
        &= 4X^3 + 6X^2 + 4X + 1\\
        f(X^3) &= (X+1)^3 - X^3\\
        &= 3X^2 + 3X + 1\\
        f(X^2) &= (X+1)^2 - X^2\\
        &= 2X + 1
    \end{align*}

    \begin{align*}
        f(X^4) - 2f(X^3) + f(X^2) &= (4X^3 + 6X^2 + 4X + 1) - 2(3X^2 + 3X + 1) + (2X + 1)\\
        &= 4X^3 + 6X^2 + 4X + 1 - 6X^2 - 6X - 2 + 2X + 1\\
        &= 4X^3 + 0X^2 + 0X + 0\\
        &= 4X^3
    \end{align*}

    Donc $X^3 = \frac{1}{4}(f(X^4) - 2f(X^3) + f(X^2))$.

    Par linéarité de $f$, on a :
    $$X^3 = f(\underbrace{\frac14 X^2(X-1)^2}_{=P})$$
\end{application}
\todo{detail}

\section{Rang d'une application linéaire}

\begin{definition}
    Soient $E,F$ deux $\Kset$-ev. et $f \in \mathcal L(E,F)$.

    Si $\Imf(f)$ est de dimension finie, alors on appelle rang de $f$ la dimension de $\Imf(f)$, et on note $\rang(f) = \underset\Kset\dim(\Imf(f))$.
\end{definition}

\begin{remark}
    \begin{enumerate}
        \item Si $F$ est de dimension finie, alors la dimension de $\Imf(f)$ est forcément finie, et par conséquent, le rang de $f$ existe et $\rang(f) \le \underset\Kset\dim(F)$. 
        \item Si $E$ est de dimension finie, alors $\rang(f)$ existe.
        \item Si $\underset\Kset\dim(E) = n \in \N^*$ et $\underset\Kset\dim(F) = p \in \N^*$, alors $\rang(f) \le \min(n,p)$.
    \end{enumerate}
\end{remark}

\begin{proof}[2e remarque précédente]
    \begin{outline}
        \1 On pose $B = (e_1,\dots,e_n)$ une base de $E$ avec $n = \dim(E)$.
        \1 Donc $\Imf(f) = f(E) = f(\Vect(e_1,\dots,e_n))$ et $\Imf(f) = \Vect(f(e_1),\dots,f(e_n))$.
        \1 Donc $\Imf(f)$ est de dimension finie ($< \infty$), et par conséquent, le rang de $f$ existe.
    \end{outline}
\end{proof}

\begin{methodbox}
    La deuxième remarque est particulièrement utile dans la pratique pour prouver que le rang d'une application linéaire existe.
\end{methodbox}

\begin{example}
    Soit $f : \begin{matrix}
    \Kset_2[X] &\to &\Kset_2[X]\\
    P &\mapsto XP' - P
    \end{matrix}$.

    Pour déterminer le rang de $f$, on peut utiliser la méthode suivante :
    \begin{outline}
        \1 On sait que $(1,X,X^2)$ est une base de $\Kset_2[X]$.
        \1 Le rang de $f$ est $\rang(f) = \rang(f(1),f(X),f(X^2))$.
        \1 On a $f(1) = -1$, $f(X) = 0$ et $f(X^2) = 2X^2$.
        \1 Donc $\rang(f) = \rang(-1,0,2X^2) = \rang(-1,X^2) = 2$.
        \1 D'où $\rang(f) = 2$.
        \1 \textit{Note additionnelle :}
        \2 On voit aussi que $\Imf(f)$ est un hyperplan\index{hyperplan} de $\Kset_2[X]$ (car il est de dimension 2 dans un espace de dimension 3).
        \2 i.e. $\exists Q \in \Kset_2[X]\setminus\Imf(f)$ tel que $\Kset_2[X] = \Imf(f) \oplus \Kset\cdot Q$.
        \2 Par exemple, on peut prendre $Q = X$ (car $X \notin \Imf(f)$), et alors $\Kset_2[X] = \Imf(f) \oplus \Kset\cdot X$.
        \2 C'est un résultat assez intéressant parce qu'il nous permet de trouver une nouvelle base de $\Kset_2[X]$ à partir d'une base de $\Imf(f)$, et d'un vecteur $Q$ qui n'est pas dans $\Imf(f)$.
    \end{outline}
\end{example}

\begin{property}
    Soient $E,F$ deux $\Kset$-ev. et $f \in \mathcal L(E,F)$.

    Si $H$ est un supplémentaire de $\Ker(f)$ dans $E$, alors $H$ est isomorphe à $\Imf(f)$, donc :
    $$H \sim \Imf(f)$$
\end{property}

\begin{proof}
    Soit $H$ un supplémentaire de $\Ker(f)$ dans $E$. Alors $E = \Ker(f) \oplus H$.

    On cherche à déterminer une application linéaire bijective de $H$ vers $\Imf(f)$.

    On pose $g : \begin{matrix}
    H &\to &\Imf(f)\\
    x &\mapsto &f(x)
    \end{matrix}$

    \textit{Au final, cela revient à considérer la restriction de $f$ à $H$.}

    \begin{outline}
        \1 $\forall x \in H$, $g(x) = f(x) \in \Imf(f)$, donc $g$ est une application de $H$ vers $\Imf(f)$.
        \1 Comme $f$ est linéaire, alors $g$ est linéaire $\in \mathcal L(H, \Imf(f))$.
        \1 Soit $x \in \Ker(g)$. Cela équivaut à dire que $x\in H et g(x) = f(x) = 0_F$, c'est-à-dire que $x \in \Ker(f)$, ce qui équivaut à dire que $H \cap \Ker(f) = \{0_E\}$, alors $x = 0_E$. Donc $g$ est injective.
        \1 On a $E = H \oplus \Ker(f)$, donc $f(E) = f(H) + f(\Ker(f)) = f(H) + \{0_F\} = f(H)$, donc $g$ est surjective, car $g(H) = f(H) = f(E) = \Imf(f)$.
        \1 Donc $g$ est bijective, et par conséquent, $H$ est isomorphe à $\Imf(f)$.
    \end{outline}
\end{proof}

\begin{theorem}[Théorème du rang]
    Soient $E,F$ deux $\Kset$-ev. et $f \in \mathcal L(E,F)$, et $\underset\Kset\dim(E) = n \in \N^* < \infty$.
    
    Alors $\underset\Kset\dim(\Ker(f)) < \infty$, et $\rang(f) = n - \underset\Kset\dim(\Ker(f))$.
\end{theorem}

\begin{proof}
    \begin{outline}
        \1 $\Ker(f)$ est un sev de $E$, donc $\Ker(f)$ est de dimension finie, et $H \subset E$ est un supplémentaire de $\Ker(f)$ dans $E$, alors $E = \Ker(f) \oplus H$.
        \1 D'après la propriété précédente, $H$ est isomorphe à $\Imf(f)$, donc $\underset\Kset\dim(H) = \underset\Kset\dim(\Imf(f)) = \rang(f)$.
        \1 Or on sait que $E = \Ker(f) \oplus H$, donc $\underset\Kset\dim(E) = \underset\Kset\dim(\Ker(f)) + \underset\Kset\dim(H)$, et par conséquent, $\rang(f) = \underset\Kset\dim(H) = n - \underset\Kset\dim(\Ker(f))$.
    \end{outline}
\end{proof}

\begin{theorem}
    Si $f \in \mathcal L(E,F)$ avec $\underset\Kset\dim(E) = \underset\Kset\dim(F) = n < \infty$, alors :
    $$f \text{ est injective} \iff f \textrm{ est surjective} \iff f \in \Isom(E,F)$$
\end{theorem}

\begin{methodbox}
    Ce théorème est particulièrement utile pour démontrer qu'une application linéaire injective liant deux ensembles de même dimension est un isomorphisme, par exemple une application injective de $\Kset_n[X]$ vers $\Kset^{n+1}$ est un isomorphisme...
\end{methodbox}

\begin{proof}
    \begin{align*}
    f \text{ est injective} &\iff \Ker(f) = \{0_E\}\\
    &\iff \rang(f) = n - \underset\Kset\dim(\Ker(f)) = n - 0 = n\\
    &\iff \dim(\Imf(f)) = n = \dim(F)\\
    &\iff \Imf(f) = F\\
    &\iff f \text{ est surjective}
    \end{align*}
\end{proof}

\begin{example}
    Soit $\Delta : \begin{matrix}
    \Kset_n[X] &\to &\Kset_n[X]\\
    P &\mapsto &P(X+1) - P(X)
    \end{matrix}$

    \begin{outline}
        \1 $\Delta$ est linéaire $\in \mathcal L(\Kset_n[X])$.
        \1 Soit $P \in \Ker(\Delta)$. Cela signifie que $\Delta(P) = 0$, donc $P(X+1) - P(X) = 0$.
        \1 Donc $\forall k \in \N, P(k) = P(0)$, et par conséquent, $P$ est constant, et $P \in \Kset$.
        \1 Donc $\Ker(\Delta) = \Kset_0[X]$.
        \1 Donc $\underset\Kset\dim(\Ker(\Delta)) = 1$.
        \1 Donc $\forall P \in \Kset_n[X]$, on a $\deg(P(X+1) - P(X)) \le n-1$, et par conséquent, $\Imf(\Delta) \subset \Kset_{n-1}[X]$.
        \1 Or  $\rang(\Delta) = n+1 - \underset\Kset\dim(\Ker(\Delta)) = n+1 - 1 = n$.
        \1 D'où $\Imf(\Delta) = \Kset_{n-1}[X]$ et $\rang(\Delta) = n$.
    \end{outline}
\end{example}

\begin{attention}
    Dans l'exemple précédent, on n'a pas $\Kset_n[X] = \Imf(\Delta) \oplus \Ker(\Delta)$, car $\Imf(\Delta) \cap \Ker(\Delta) = \{0_{\Kset_n[X]}\}$.

    En général, la somme directe n'est pas vérifiée !
\end{attention}

\begin{consequence}
    Soit $f \in \mathcal L(E)$ avec $\underset\Kset\dim(E) = n < \infty$. Alors $f$ est injective $\iff$ $f$ est surjective $\iff$ $f \in \GL(E)$.
\end{consequence}

\begin{property}[Rang d'une composée d'applications linéaires]
    Soient $E, F, G$ trois $\Kset$-ev. et $f \in \mathcal L(E,F)$, $g \in \mathcal L(F,G)$, tels que $\underset\Kset\dim(E) < \infty$ et $\underset\Kset\dim(F) < \infty$.

    \begin{enumerate}
        \item $\rang(g\circ f) = \rang(f) - \underset\Kset\dim(\underbrace{\Ker(g) \cap \Imf(f)}_{=\Ker(g\circ f)})$.
        \item Si $g$ est injective, alors $\rang(g\circ f) = \rang(f)$.
        \item Si $f$ est surjective, alors $\rang(g\circ f) = \rang(g)$.
        \item \textbf{Les applications linéaires bijectives ne changent pas le rang :} Si $f \in \Isom(E,F)$ et $g \in \Isom(F,G)$, alors $\rang(g\circ f) = \rang(f) = \rang(g)$.
    \end{enumerate}
\end{property}

\begin{proof}
    \begin{outline}[enumerate]
    \1 On peut considérer $g\circ f$ comme l'application linéaire partant de $\Imf(f)$ vers $G$, et $\Imf(f)$ est de dimension finie, donc le rang de $g\circ f$ existe, et on applique donc le théorème du rang à $g\circ f$ pour retrouver le résultat.
    \1 Si $g$ est injective, alors $\Ker(g) = \{0_F\}$, donc $\Ker(g) \cap \Imf(f) = \{0_F\}$, et par conséquent, $\rang(g\circ f) = \rang(f)$.
    \1 Si $f$ est surjective, alors $\Imf(f) = F$, donc $\Ker(g) \cap \Imf(f) = \Ker(g)$, et par conséquent, $\rang(g\circ f) = \rang(f) - \underset\Kset\dim(\Ker(g)) = \underset\Kset\dim(F) - \underset\Kset\dim(\Ker(g)) = \rang(g)$.
    \end{outline}
\end{proof}

\chapter{Projections et symétries}

\begin{property}
    Soient $E$ un $\Kset$-ev. et $(E_i)_{1\le i \le n}$ une famille de sev de $E$ tels que $E = \bigoplus_{i=1}^n E_i$.

    Soient $u_i \in \mathcal L(E_i,F)$ pour $i=1,\dots,n$.

    Alors $\exists! u \in \mathcal L(E,F)$ tel que $U_{|E_i} = u_i$ pour $i=1,\dots,n$ (on note $u_{|E_i}$ la restriction de $u$ à $E_i$).

    $u$ est donnée par $u : \begin{matrix}
    E &\to &F\\
    \sum_{i=1}^n x_i &\mapsto &\sum_{i=1}^n u_i(x_i)
    \end{matrix}$ avec $x_i \in E_i$ pour $i=1,\dots,n$.
\end{property}

\begin{proof}
    À faire en tant qu'exercice

    (il faut montrer que $u$ est bien définie, linéaire, que sa restriction à $E_i$ est bien égale à $u_i$, et que $u$ est la seule application linéaire vérifiant ces propriétés).
\end{proof}

\section{Projections et projecteurs}

\begin{definition}[Projection]
    Soient $E$ un $\Kset$-ev. et $F,G$ deux sous-espaces vectoriels de $E$ tels que $E = F \oplus G$.

    On appelle projection de $E$ sur $F$ parallèlement à $G$ l'endomorphisme $p$ de $E$ vérifiant :
    $$\begin{cases}
    \forall x \in F, p(x) = x\\
    \forall x \in G, p(x) = 0_E
    \end{cases}$$
\end{definition}

\begin{remark}
    On sait que $\forall x \in E, \exists! (x_1,x_2) \in F \times G$ tel que $x = x_1 + x_2$, et alors $p(x) = p(x_1 + x_2) = p(x_1) + p(x_2) = x_1 + 0_E = x_1$.
\end{remark}

\begin{definition}[Projecteur]
    Soit $f \in \mathcal L(E)$. $f$ est dit projecteur de $E$ si $f^2 = f$.
\end{definition}

\begin{remark}
    $f$ est une projection sur $F$ parallèlement à $G$ implique que $f$ est un projecteur, mais la réciproque n'est pas forcément vérifiée.
\end{remark}

\begin{proof}
    Comme on sait que $p(x) = x_1$ pour $x = x_1 + x_2$ avec $x_1 \in F$ et $x_2 \in G$, alors $p^2(x) = p(p(x)) = p(x_1) = x_1 = p(x)$, et par conséquent, $p^2 = p$, donc $p$ est un projecteur.
\end{proof}

\begin{consequence}
    Si $f$ est un projecteur, alors $\forall k \ge 1, f^k = f$.
\end{consequence}

% \begin{exercise}
%     Montrer que:
%     \begin{itemize}
%         \item $f \in \mathcal L(E)$ est un projecteur $\implies$ $E = \Imf(f) \oplus \Ker(f)$.
%         \item $f \in \mathcal L(E)$ est un projecteur $\implies$ $f$ est une projection sur $\Imf(f)$ parallèlement à $\Ker(f)$.  
%     \end{itemize}
% \end{exercise}

\begin{property}
    Soit $f \in \mathcal L(E)$.
    
    Alors $f$ est un projecteur $\implies$ $E = \Imf(f) \oplus \Ker(f)$.
\end{property}

\begin{proof}
    Réalisée en tant qu'exercice.

    \begin{outline}
        \1 On montre que l'intersection de $\Imf(f)$ et $\Ker(f)$ est réduite au vecteur nul.
        \2 On prend $x \in \Imf(f) \cap \Ker(f)$, alors $\exists y \in E$ tel que $x = f(y)$, et $f(x) = 0_E$, donc $f(f(y)) = 0_E$, donc $f(y) = 0_E$, donc $x = f(y) = 0_E$.
        \1 On montre que $E = \Imf(f) + \Ker(f)$.
        \2 On remarque directement que $\forall x \in E, f^2(x) = f(x) \implies f(f(x)) = f(x) \implies f(f(x)-x)=0_E \implies f(x)-x \in \Ker(f)$. Donc $x = f(x) + (x - f(x))$ avec $f(x) \in \Imf(f)$ et $x - f(x) \in \Ker(f)$, et par conséquent, $E = \Imf(f) + \Ker(f)$.
        \1 Donc $E = \Imf(f) \oplus \Ker(f)$.
    \end{outline}
\end{proof}

\begin{remark}
    Si $f$ est un projecteur, alors $f$ est une projection sur $\Imf(f)$ parallèlement à $\Ker(f)$.
\end{remark}

\begin{proof}
    On a $\forall x \in \Ker(f), f(x) = 0_E$ (par définition) et $\forall x \in \Imf(f), \exists y \in E$ tel que $x = f(y)$, donc $f(x) = f(f(y)) = f(y) = x$ (car $f$ est un projecteur), et par conséquent, sachant que $E = \Imf(f) \oplus \Ker(f)$, $f$ est une projection sur $\Imf(f)$ parallèlement à $\Ker(f)$.
\end{proof}

\begin{remark}
    Si $f$ est un projecteur, alors $\Imf(f) = \Ker(f - \Id_E)$.
\end{remark}

\begin{proof}
    La démonstration se fait en utilisant le lemme des noyaux, qui sera vu prochainement. On peut passer de $f^2 = f$ à un polynôme annulateur de $f$ qui est $X^2 - X$, et alors on peut factoriser ce polynôme annulateur en $(X)(X-1)$, et alors on peut appliquer le lemme des noyaux pour trouver que $E = \Ker(f) \oplus \Ker(f - \Id_E)$, et par conséquent, $\Imf(f) = \Ker(f - \Id_E)$.
\end{proof}

\begin{remark}[Bijectivité d'une projection]
    Si $f$ est un projecteur, alors $f$ est bijective $\iff$ $f = \Id_E$. Pour que $f$ soit bijective, il faut que $\Ker(f) = \{0_E\}$, et alors $E = \Imf(f)$, et par conséquent, $f$ est une projection sur $E$ parallèlement à $\{0_E\}$, et alors $f = \Id_E$ ($f(x) = x$ pour tout $x \in E$).
    
\end{remark}

\section{Symétrie et involution}

\begin{definition}[Symétrie]
    Soient $E$ un $\Kset$-ev. et $F,G$ deux sous-espaces vectoriels de $E$ tels que $E = F \oplus G$.

    On appelle symétrie par rapport à $F$ parallèlement à $G$ l'endomorphisme $s$ de $E$ vérifiant :
    $$\begin{cases}
    \forall x \in F, s(x) = x\\
    \forall x \in G, s(x) = -x
    \end{cases}$$
\end{definition}

\begin{remark}
    On sait que $\forall x \in E, \exists! (x_1,x_2) \in F \times G$ tel que $x = x_1 + x_2$, et alors $s(x) = s(x_1 + x_2) = s(x_1) + s(x_2) = x_1 - x_2$.

    Donc $s^2(x) = s(x_1 - x_2) = s(x_1) - s(x_2) = x_1 + x_2 = x$, et par conséquent, $\boxed{s^2 = \Id_E}$.
\end{remark}

\begin{remark}
    Soit $p$ la projection de $E$ sur $F$ parallèlement à $G$, et soit $q$ la projection de $E$ sur $G$ parallèlement à $F$.

    On a donc $p+q = \Id_E$ (car $p(x) + q(x) = x_1 + x_2 = x$ pour tout $x \in E$), et $\boxed{s = p - q}$ (car $p(x) - q(x) = x_1 - x_2 = s(x)$ pour tout $x \in E$).
\end{remark}

\begin{remark}
    On a aussi $\Ker(s-\Id_E) = F$ et $\Ker(s+\Id_E) = G$ (car $s(x) = x$ pour tout $x \in F$, et $s(x) = -x$ pour tout $x \in G$).
\end{remark}

\begin{definition}[Involution]
    Soit $f \in \mathcal L(E)$. $f$ est dit involution de $E$ si $f^2 = \Id_E$.
\end{definition}

\begin{remark}
    Toute symétrie est une involution.
\end{remark}

\begin{property}
    Soit $f \in \mathcal L(E)$. Alors $f$ est une involution $\implies$ $E = \Ker(f - \Id_E) \oplus \Ker(f + \Id_E)$.
\end{property}

\begin{proof}
    \begin{outline}
        \1 Montrons que $\Ker(f - \Id_E) \cap \Ker(f + \Id_E) = \{0_E\}$.
        \2 Soit $x \in \Ker(f - \Id_E) \cap \Ker(f + \Id_E)$, alors $f(x) = x$ et $f(x) = -x$, donc $x = 0_E$, et par conséquent, $\Ker(f - \Id_E) \cap \Ker(f + \Id_E) = \{0_E\}$.
        \1 Montrons que $E = \Ker(f - \Id_E) + \Ker(f + \Id_E)$.
        \2 Soit $x \in E$. On va procéder par analyse-synthèse.
        \2 \textit{Analyse :}
        \3 On suppose que $x = x_1 + x_2$ avec $x_1 \in \Ker(f - \Id_E)$ et $x_2 \in \Ker(f + \Id_E)$.
        \3 Alors $f(x) = f(x_1) + f(x_2) = x_1 - x_2$.
        \3 Donc $\begin{cases}x_1 + x_2 = x\\ x_1 - x_2 = f(x)\end{cases}$, et par conséquent, $\begin{cases}x_1 = \frac{x + f(x)}{2}\\ x_2 = \frac{x - f(x)}{2}\end{cases}$.
        \2 \textit{Synthèse :}
        \3 On pose $x_1 = \frac{x + f(x)}{2}$ et $x_2 = \frac{x - f(x)}{2}$.
        \3 Alors $f(x_1) = \frac{f(x) + f^2(x)}{2} = \frac{f(x) + x}{2} = x_1$.
        \3 Et $f(x_2) = \frac{f(x) - f^2(x)}{2} = \frac{f(x) - x}{2} = -x_2$.
        \3 Donc $x_1 \in \Ker(f - \Id_E)$ et $x_2 \in \Ker(f + \Id_E)$.
        \3 Alors $x = x_1 + x_2$.
        \1 Donc $E = \Ker(f - \Id_E) \oplus \Ker(f + \Id_E)$.
    \end{outline}
\end{proof}

\begin{remark}
    Si $f$ est une involution, alors $f$ est une symétrie par rapport à $\Ker(f - \Id_E)$ parallèlement à $\Ker(f + \Id_E)$.
\end{remark}

\begin{theorem}[Lemme]
    Soient $E$ un $\Kset$-ev. et $f,g \in \mathcal L(E)$.
    $$f \circ g = 0 \iff \Imf(g) \subset \Ker(f)$$
\end{theorem}

\begin{proof}
    \begin{outline}
        \1 ($\implies$) Soit $y \in \Imf(g)$, alors $\exists x \in E$ tel que $y = g(x)$, et alors $f(y) = f(g(x)) = 0_E$, et par conséquent, $y \in \Ker(f)$, et par conséquent, $\Imf(g) \subset \Ker(f)$.
        \1 ($\impliedby$) Soit $x \in E$, alors $g(x) \in \Imf(g)$, et par conséquent, $g(x) \in \Ker(f)$, donc $f(g(x)) = 0_E$, et par conséquent, $f\circ g = 0$.
    \end{outline}
\end{proof}

\begin{theorem}[Lemme de la décomposition des noyaux]
    Soient $E$ un $\Kset$-ev. et $f \in \mathcal L(E)$ et $P \in \Kset[X]$ tel que $P = P_1\cdot P_2$ avec $P$ un polynôme annulateur de $f$ et $\pgcd(P_1,P_2) = 1$.
    $$E = \Ker(P_1(f)) \oplus \Ker(P_2(f))$$
\end{theorem}

\begin{proof}
    \begin{outline}
        \1 En utilisant le théorème de Bézout, on trouve $U,V \in \Kset[X]$ tels que $UP_1 + VP_2 = 1$.
        \1 On passe au polynôme d'endomorphisme pour trouver $(UP_1)(f) + (VP_2)(f) = \Id_E$.
        \1 Cela nous donne :
        $$\begin{cases}
        U(f)\circ P_1(f) + V(f)\circ P_2(f) = \Id_E\\
        P_1(f)\circ U(f) + P_2(f)\circ V(f) = \Id_E
        \end{cases}$$
        \1 Soit $x \in \Ker(P_1(f)) \cap \Ker(P_2(f))$.
        \1 Donc $x \in E$ et $\begin{cases} P_1(f)(x) = 0_E\\ P_2(f)(x) = 0_E\end{cases}$.
        \1 Donc 
    \end{outline}
\end{proof}
\todo{recop}

\begin{remark}[Généralisation du lemme de la décomposition des noyaux]
    Si $P$ n'est pas forcément annulateur de $f$, alors :
    $$\Ker(P(f)) = \Ker(P_1(f)) \oplus \Ker(P_2(f))$$

    Si $P(f) = 0$, alors :
    $$\Ker(0_{\mathcal L(E)}) = E = \Ker(P_1(f)) \oplus \Ker(P_2(f))$$
\end{remark}

\begin{remark}
    Par récurrence, si $P = \prod_{i=1}^{r}P_k$ avec $P_1,\dots,P_r$ des polynômes annulateurs de $f$ tels que $\pgcd(P_i,P_j) = 1$ pour $i \neq j$, alors :
$$E = \bigoplus_{i=1}^{r} \Ker(P_i(f))$$
\end{remark}

\begin{exercise}
    On pose $f : \begin{matrix}
        R^2 &\to &R^2\\
        (x,y) &\mapsto &(4x-y, 2x+y)
    \end{matrix}$
    \begin{enumerate}
        \item Montrer que $X^2-5X+6$ est un polynôme annulateur de $f$.
        \item Calculer $f^n$ pour $n \in \N^*$.
    \end{enumerate}

    \textbf{Solution :}

    \begin{outline}
        \1 \textbf{1.} Simple.
        \1 \textbf{2.} On effectue la division euclidienne de $X^n$ par $P$. 
        \2 On trouve $X^n = Q_n(X)(X^2 - 5X + 6) + a_n X + b_n$ avec $Q_n \in \R_n[X]$ et $a_n,b_n \in \R$.
        \2 Pour $X=2$, on a $2^n = a_n \cdot 2 + b_n$.
        \2 Pour $X=3$, on a $3^n = a_n \cdot 3 + b_n$.
        \2 En résolvant ce système de deux équations à deux inconnues, on trouve $a_n = 3^n - 2^n$ et $b_n = 2 \cdot 3^n - 3 \cdot 2^n$.
        \2 Donc $X^n = Q_n(X)(X^2 - 5X + 6) + (3^n - 2^n)X + (2 \cdot 3^n - 3 \cdot 2^n)$.
        \1 Donc $f^n = (3^n - 2^n)f + (2 \cdot 3^n - 3 \cdot 2^n)\Id_{R^2}$, sachant que $P(f) = 0$ donc le terme $Q_n(f)(f^2 - 5f + 6\Id_{R^2})$ est égal à $0$.
        \1 On a aussi le résultat suivant :
        $$R^2 = \Ker(f - 2\Id_{R^2}) \oplus \Ker(f - 3\Id_{R^2})$$
        (en décomposant $P$, et sachant que $(X-2) \wedge (X-3) = 1$).
    \end{outline}

    Une méthode alternative consiste à procéder par récurrence à partir de l'expression $f^2 = 5f - 6\Id_{R^2}$, et à trouver une relation de récurrence pour $a_n$ et $b_n$ à partir de l'expression $f^n = a_n f + b_n \Id_{R^2}$. On trouvera lors de la récurrence $\begin{cases} a_{n+1} = 5a_n + b_n\\ b_{n+1} = -6a_n\end{cases}$, et alors on peut résoudre ce système de récurrence pour trouver les expressions de $a_n$ et $b_n$ (qui seront les mêmes que celles trouvées précédemment).
\end{exercise}

\chapter{Formes linéaires}

\begin{definition}
    Soit $E$ un $\Kset$-ev.

    Toute application linéaire $\varphi$ de $E$ vers $\Kset$ est appelée forme linéaire de $E$.

    On note $E^*$ l'ensemble des formes linéaires, i.e. $E^* = \mathcal L(E,\Kset)$, et il est appelé \textit{espace dual} de $E$.
\end{definition}

\begin{attention}
    Il ne faut pas confondre la notation $E^*$ pour l'espace dual de $E$ avec la notation $E^*$ pour l'ensemble $E$ privé du vecteur nul (i.e. $E^* = E \setminus \{0_E\}$) (ou des éléments inversibles...).
\end{attention}

\begin{example}
    L'application $\varphi : \begin{matrix}
    \mathcal C([a,b],\Kset) \to \Kset\\
    f \mapsto \int_a^b f(t) dt
    \end{matrix}$ avec $a < b$ est une forme linéaire de $\mathcal C([a,b],\Kset)$, et par conséquent, $\varphi \in \mathcal C([a,b],\Kset)^*$.
\end{example}

\begin{property}
    Soient $E$ un $\Kset$-ev. et $H$ un sev. de $E$.

    $H$ est un hyperplan de $E$ $\iff$ $\exists \varphi \in E^* \setminus \{0_{E^*}\}$ telle que $H = \Ker(\varphi)$.
    
    i.e. : Tout hyperplan $H$ de $E$ est le noyau d'une forme linéaire non nulle de $E$.
\end{property}

% kholle spetsialne

\begin{proof}
    $\Longrightarrow$ Soit $H$ un hyperplan de $E$.
Il existe $a \in E \setminus \{0_E\}$ tel que $E = H \oplus \mathbb{K} \cdot a$. On peut exprimer tout vecteur $x \in E$ sous la forme 

$$\forall x \in E, \exists! (h_x, \lambda_x) \in H \times \mathbb{K} \text{tel que } x = h_x + \lambda_x \cdot a$$

On pose $\varphi : E \to \mathbb{K}$ l'application qui à $x = h_x + \lambda_x \cdot a$ associe $\varphi(x) = \lambda_x$.
D'où $\forall x \in E, \exists! h_x \in H$ tel que $x = h_x + \varphi(x) \cdot a$.

\begin{itemize}
    \item Montrons que $\varphi$ est linéaire : soient $x, y \in E$ et $\alpha \in \mathbb{K}$.
    On a $x = h_x + \varphi(x) \cdot a$ et $y = h_y + \varphi(y) \cdot a$.
    D'où $x + \alpha y = h_x + \alpha h_y + (\varphi(x) + \alpha \varphi(y)) \cdot a$.
    Or, par définition, on a aussi $x + \alpha y = h_{x+\alpha y} + \varphi(x + \alpha y) \cdot a$.
    Donc par unicité de la décomposition, on obtient :
    $\varphi(x + \alpha y) = \varphi(x) + \alpha \varphi(y)$, soit $\varphi \in \mathcal{L}(E, \mathbb{K})$.
    
    \item Noyau de $\varphi$ :
    $x \in \ker(\varphi) \iff \varphi(x) = 0 \iff x = h_x \in H$, donc $\ker(\varphi) = H$.
    
    \item Non-nullité de $\varphi$ :
    Comme $a \in E$, on a $a = 0_E + 1 \cdot a = h_a + \varphi(a) \cdot a$.
    D'où par unicité, $\varphi(a) = 1 \neq 0$.
    \end{itemize}
    Conclusion du sens direct : $\exists \varphi \in \mathcal{L}(E, \mathbb{K}) \setminus \{0\}$ tel que $\ker(\varphi) = H$.

    $\Longleftarrow$ Supposons qu'il existe $\varphi \in \mathcal{L}(E, \mathbb{K}) \setminus \{0\}$ tel que $\ker(\varphi) = H$.
    Montrons que $H = \ker(\varphi)$ est un hyperplan de $E$, c'est-à-dire que $\forall b \in E \setminus H, E = H \oplus \mathbb{K} \cdot b$.
    Soit $b \in E \setminus H$.

    \begin{itemize}
        \item Intersection : soit $x \in H \cap (\mathbb{K} \cdot b)$.
        On a $\varphi(x) = 0$ et $\exists \lambda \in \mathbb{K}$ tel que $x = \lambda \cdot b$.
        D'où $\varphi(x) = 0 = \varphi(\lambda \cdot b) = \lambda \cdot \varphi(b)$.
        Si $\lambda \neq 0$, alors $\varphi(b) = 0 \implies b \in \ker(\varphi) = H$, ce qui est absurde car $b \notin H$.
        Donc $\lambda = 0 \implies x = 0_E$, d'où $H \cap (\mathbb{K} \cdot b) = \{0_E\}$.
        
        \item Somme : soit $x \in E$, cherchons $\lambda \in \mathbb{K}$ tel que $x - \lambda \cdot b \in H$.
        $x - \lambda \cdot b \in H \iff \varphi(x - \lambda \cdot b) = 0 \iff \varphi(x) - \lambda \cdot \varphi(b) = 0 \iff \lambda = \frac{\varphi(x)}{\varphi(b)}$ (car $\varphi(b) \neq 0$ puisque $b \notin H$).
        On peut alors écrire :
        $$x = \underbrace{x - \frac{\varphi(x)}{\varphi(b)} \cdot b}_{\in H} + \underbrace{\frac{\varphi(x)}{\varphi(b)} \cdot b}_{\in \mathbb{K} \cdot b}$$
    \end{itemize}
    D'où $E = H \oplus \mathbb{K} \cdot b$, ce qui montre que $H$ est un hyperplan de $E$.
\end{proof}

\begin{exercise}
    On pose $H = \{(x,y,z) \in \R^3 | x -y + z = 0\}$ et $D = \R\cdot(1,1,2)$.
    \begin{enumerate}
        \item Montrer que $R^3 = H \oplus D$.
        \item Donner l'expression de la projection de $H$ parallèlement à $D$.
    \end{enumerate}

    \textbf{Solution :}

    \begin{outline}
        \1 La première question est simple. On vérifie que $\underset{\R^3}\dim(H) = 2$ et que $(1,1,2) \notin H$, donc $R^3 = H \oplus D$.
        \1 Sinon, on peut aussi poser $\varphi : (x,y,z) \mapsto x - y + z$, alors $\varphi$ est une forme linéaire de $\R^3$, et $H = \Ker(\varphi)$, et alors $D$ est un supplémentaire de $H$ dans $\R^3$.
        \1 Soit $(x,y,z) \in \R^3$.
        \2 
    \end{outline}
\end{exercise}
\todo{recop}

\chapter{Compléments}

\begin{definition}[Valeur propre, vecteur propre]
    Soient $E$ un $\Kset$-ev, $f \in \mathcal L(E)$ et $\lambda \in \Kset$.

    $\lambda$ est dite \textbf{valeur propre} de $f$ s'il existe un vecteur non nul $u \in E\setminus\{0_E\}$ tel que $f(u) = \lambda u$.

    Le vecteur $u$ est alors appelé \textbf{vecteur propre} de $f$ associé à la valeur propre $\lambda$.
\end{definition}

\begin{definition}[Spectre]
    On appelle l'ensemble des valeurs propres de $f$ le spectre de $f$, et on le note $\Sp_\Kset(f)$.
\end{definition}

\begin{definition}[Sous-espace propre]
    Si $\lambda$, une valeur propre de $f$, appartient à $\Sp_\Kset(f)$, alors $\Ker(f-\lambda \Id_E)$ est appelé le \textbf{sous-espace propre} de $f$ associé à la valeur propre $\lambda$, et on le note $E_\lambda = \Ker(f - \lambda \Id_E)$.
\end{definition}

\begin{remark}
    \begin{enumerate}
        \item $\lambda \in \Sp_\Kset(f) \iff \Ker(f - \lambda \Id_E) \neq \{0_E\}$.
        \item Si $\underset\Kset\dim(E) < +\infty$, alors $\lambda \in \Sp_\Kset(f) \iff f - \lambda \Id_E$ n'est pas bijective.
    \end{enumerate}
\end{remark}

\begin{example}
    \begin{outline}
        $f : \begin{matrix}
        \R^2 &\to &\R^2\\
        (x,y) &\mapsto &(4x-y, 2x+y)
        \end{matrix}$

        \1 Soit $\lambda \in \Sp_\R(f)$.
        
        \1 Cela équivaut à dire que $\exists (u,v) \in \R^2 \setminus \{(0,0)\}$ tel que $f(u,v) = \lambda (u,v)$, c'est à dire que $\begin{cases}4u - v = \lambda u\\ 2u + v = \lambda v\end{cases}$, et alors $\begin{cases} (4-\lambda)u - v = 0\\ 2u + (1-\lambda)v = 0\end{cases}$.

        \1 Si $u = 0$, alors $v = 0$, ce qui est impossible, donc $u \neq 0$, et alors $v = (4-\lambda)u$.

        \1 On trouve alors $2u + (1-\lambda)(4-\lambda)u = 0$, et alors $2 + (1-\lambda)(4-\lambda) = 0$, et alors $\lambda^2 - 5\lambda + 6 = 0$, et alors $\lambda \in \{2,3\}$.

        \1 Donc $\Sp_\R(f) \subset \{2,3\}$ (l'équivalence n'est plus valable à un certain point dans notre raisonnement).

        \1 Soient $(a,b) \in \Ker(f - 2\Id_{\R^2})$ et $(c,d) \in \Ker(f - 3\Id_{\R^2})$.

        \begin{align*}
            \ona f(a,b) = 2(a,b) &\iff \begin{cases}4a - b = 2a\\ 2a + b = 2b\end{cases}\\
            &\iff \begin{cases} 2a - b = 0\\ 2a - b = 0\end{cases}\\
            &\iff b = 2a
        \end{align*}

        \1 i.e. $E_2 = \Ker(f - 2\Id_{\R^2}) = \R\cdot(1,2)$. N'importe quel vecteur de $E_2$ est un vecteur propre de $f$ associé à la valeur propre $2$.

        \begin{align*}
            \ona f(c,d) = 3(c,d) &\iff \begin{cases}4c - d = 3c\\ 2c + d = 3d\end{cases}\\
            &\iff \begin{cases} c - d = 0\\ 2c - 2d = 0\end{cases}\\
            &\iff d = c
        \end{align*}

        \1 i.e. $E_3 = \Ker(f - 3\Id_{\R^2}) = \R\cdot(1,1)$. N'importe quel vecteur de $E_3$ est un vecteur propre de $f$ associé à la valeur propre $3$.
    \end{outline}
\end{example}

\begin{example}
    \begin{outline}
        $f : \begin{matrix}
            \R^2 &\to &\R^2\\
            (x,y) &\mapsto &(-y,x)
        \end{matrix}$
        \1 Soit $\lambda \in \Sp_\R(f)$.
        \1 Cela équivaut à dire que $\exists (u,v) \in \R^2 \setminus \{(0,0)\}$ tel que $f(u,v) = \lambda (u,v)$, c'est à dire que $\begin{cases}-v = \lambda u\\ u = \lambda v\end{cases}$, et alors $\begin{cases} -v - \lambda u = 0\\ u - \lambda v = 0\end{cases}$.
        \1 Si $u = 0$, alors $v = 0$, ce qui est impossible, donc $u \neq 0$, et alors $v = -\lambda u$.
        \1 On trouve alors $u - \lambda(-\lambda u) = 0$, et alors $1 + \lambda^2 = 0$, ce qui est impossible, donc $\Sp_\R(f) = \emptyset$.
        \1 On a $f^2(x,y) = f(-y,x) = (-x,-y) = -\Id_{\R^2}(x,y)$, et alors $f^2 + \Id_{\R^2} = 0$, donc $P(X) = X^2 + 1$ est un polynôme annulateur de $f$. $P$ n'est pas scindé sur $\R$.\index{Polynôme annulateur non scindé dans $\R$}
        Si on travaille dans $\C$, alors $\Sp_\C(f) = \{i,-i\}$.
    \end{outline}
\end{example}

\begin{definition}[Algèbre]
    Soit $A$ un ensemble non-vide muni de deux LCI $+$ et $\times$ et d'une LCE $\cdot$.
    $(A,+,\times,\cdot)$ est dit \textbf{$\Kset$-algèbre} si :
    \begin{enumerate}
        \item $(A,+,\times)$ est un anneau.
        \item $(A,+,\cdot)$ est un $\Kset$-ev.
        \item $\forall \lambda \in \Kset, \forall x,y \in A, \begin{cases} \lambda(x\times y) = (\lambda x)\times y = x \times (\lambda y) \\ (\alpha\lambda)x = \alpha(\lambda x) = \lambda(\alpha x)\end{cases}$.
    \end{enumerate}

    Si, de plus, la loi $\times$ est commutative dans $A$, alors on dit que $A$ est une $\Kset$-algèbre commutative.
\end{definition}

\begin{example}
    \begin{itemize}
        \item $\C$ est une $\R$-algèbre commutative.
        \item $(\Kset[X],+,\times,\cdot)$ est une $\Kset$-algèbre commutative.
        \item $(\Kset^\R, +, \times, \cdot)$ est une $\Kset$-algèbre commutative.
        \item Si $E$ est un $\Kset$-ev., alors $(\mathcal L(E), +, \circ, \cdot)$ est une $\Kset$-algèbre (non commutative).
    \end{itemize}
\end{example}

% je suis parti voir les cours de c2026 pour rien parce qu'on a pas écrit cette définition au final =)

\begin{definition}[Sous-algèbre]
    Soit $A$ une $\Kset$-algèbre et $B \subset A$.
    $B$ est une \textbf{sous-algèbre} de $A$ si :
    \begin{enumerate}
        \item $1_A \in B$.
        \item $\forall x,y \in B, \forall \lambda \in \Kset, \begin{cases} x + \lambda y \in B\\ x \times y \in B\end{cases}$.
    \end{enumerate}
\end{definition}

\begin{center}
    \rule{0.5\textwidth}{1pt}\\
    Fin du Chapitre XIV-B.\\
    \rule{0.5\textwidth}{1pt}
\end{center}

\printindex
\end{document}